% Part: lambda-calculus
% Chapter: introduction
% Section: beta

\documentclass[../../../include/open-logic-section]{subfiles}

\begin{document}

\olfileid{lam}{int}{bet}
\olsection{$\beta$-न्यूनीकरण}

$(\lambd[m][(\lambd[y][y]) m])$ पाहिल्यावर त्याचा
$\lambd[m][m]$ शी काही संबंध असावा असे वाटते:
दुसरे पद पहिल्याचे ``सोपीकरण'' केल्यावर मिळायला हवे.
\emph{$\beta$-न्यूनीकरण} ही कल्पना औपचारिकपणे पकडते.

\begin{defn}[$\beta$-संकुचन, $\bredone$] \ollabel{defn:betacontr}
  \emph{$\beta$-संकुचन} ($\bredone$) हा पदांवरील
  सर्वांत लहान रचनांशी सुसंगत संबंध आहे जो पुढील अट पूर्ण करतो:
  \[
    (\lambd[x][N])Q \bredone \Subst{N}{Q}{x}
  \]
  $P \bredone Q$ असेल, तर $P$ चे $Q$ मध्ये
  \emph{$\beta$-संकुचन} झाले असे म्हणतो.
  $(\lambd[x][N])Q$ या रूपाच्या पदाला \emph{रेडेक्स} म्हणतात.
\end{defn}

\begin{prob} \ollabel{prob:def}
  \olref[lam][syn][alp]{defn:aconvone} मध्ये बद्ध चराच्या
  नामांतरासाठी जशा मांडल्या, तशाच $\beta$-संकुचनाच्या
  सममूल्य आगमनाधारित व्याख्या सविस्तर लिहा.
\end{prob}

\begin{defn}[$\beta$-न्यूनीकरण, $\bred$] \ollabel{defn:betared}
  \emph{$\beta$-न्यूनीकरण} ($\bred$) हा $\bredone$
  समाविष्ट करणारा पदांवरील सर्वांत लहान परावर्ती आणि
  संक्रमक संबंध आहे. $P \bred Q$ असेल, तर $P$ चे
  $Q$ मध्ये $\beta$-न्यूनीकरण झाले असे म्हणतो.
\end{defn}

संदर्भ स्पष्ट असेल, तर $\bredone$ ऐवजी $\redone$
आणि $\bred$ ऐवजी $\red$ असे लिहू.

अनौपचारिकपणे, $M \bred N$ तेव्हाच आणि तेव्हाच
$\beta$-संकुचनाच्या शून्य किंवा अधिक पायऱ्यांनी
$M$ चे $N$ मध्ये परिवर्तन करता येते.

\begin{defn}[$\beta$-सामान्य]
ज्या पदाचे आणखी $\beta$-संकुचन करता येत नाही,
त्याला \emph{$\beta$-सामान्य} म्हणतात.
\end{defn}

$M \bred N$ आणि $N$ हे $\beta$-सामान्य असेल, तर
$N$ ला~$M$ चे \emph{सामान्य रूप} म्हणतो. एखाद्या पदाचे
सामान्य रूप एकमेव असते का, असा प्रश्न पडू शकतो.
उत्तर होय आहे, हे आपण पुढे पाहू.

काही उदाहरणे पाहू.
\begin{enumerate}
\item आपल्याला मिळते:
\begin{align*}
(\lambd[x][xxy]) \lambd[z][z] & \redone (\lambd[z][z])(\lambd[z][z]) y \\
& \redone (\lambd[z][z]) y \\
& \redone y
\end{align*}
\item पदाचे ``सोपीकरण'' केल्यावर ते प्रत्यक्षात अधिक गुंतागुंतीचेही होऊ शकते:
\begin{align*}
(\lambd[x][xxy])(\lambd[x][xxy]) & \redone (\lambd[x][xxy])(\lambd[x][xxy])y \\
& \redone (\lambd[x][xxy])(\lambd[x][xxy])yy \\
& \redone \dots
\end{align*}
\item पद तसेचही राहू शकते:
\[
(\lambd[x][xx])(\lambd[x][xx]) \redone (\lambd[x][xx])(\lambd[x][xx])
\]
\item तसेच काही पदांचे एकाहून अधिक मार्गांनी न्यूनीकरण
  करता येते. उदाहरणार्थ,
\[
(\lambd[x][(\lambd[y][yx]) z]) v \redone (\lambd[y][yv]) z
\]
हे बाहेरील उपयोजनाचे संकुचन करून मिळते; आणि
\[
(\lambd[x][(\lambd[y][yx]) z]) v \redone (\lambd[x][zx]) v
\]
हे आतील उपयोजनाचे संकुचन करून मिळते. तथापि,
या प्रसंगी दोन्ही पदांचे पुढील न्यूनीकरण एकाच
$zv$ या पदात होते.
\end{enumerate}

शेवटच्या उदाहरणातील समान परिणाम हा योगायोग नाही.
तो लॅम्डा कलनाचा चर्च–रॉसर गुणधर्म म्हणून ओळखला
जाणारा खोल आणि महत्त्वाचा गुणधर्म दाखवतो.

\begin{digress}
  सर्वसाधारणपणे, एखाद्या पदाचे $\beta$-न्यूनीकरण
  अनेक प्रकारे करता येते. म्हणून न्यूनीकरणाच्या अनेक
  पद्धती मांडल्या गेल्या आहेत. त्यांपैकी सर्वाधिक प्रचलित
  \emph{स्वाभाविक पद्धत} आहे. ती नेहमी \emph{डावीकडील
  सर्वप्रथम} रेडेक्सचे संकुचन करते; रेडेक्सची जागा
  पदात तो जिथे सुरू होतो त्यावरून ठरवली जाते.
  या पद्धतीचा उपयुक्त गुणधर्म असा: एखाद्या पदाचे
  कोणत्याही पद्धतीने सामान्य रूपात न्यूनीकरण करता येते
  तेव्हाच स्वाभाविक पद्धतीनेही ते करता येते. पुढे
  वेगळे स्पष्ट केलेले नसेल, तर आपण स्वाभाविक पद्धत वापरू.
\end{digress}

\begin{defn}[$\beta$-सममूल्यता, $\equal$]
  \emph{$\beta$-सममूल्यता} ($\equal$) हा संबंध आगमनाने
  पुढीलप्रमाणे परिभाषित करतात:
  \begin{enumerate}
  \item $M \equal M$.
  \item $M \equal N$ असेल, तर $N \equal M$.
  \item $M \equal N$, $N \equal O$ असतील, तर $M \equal O$.
  \item $M \equal N$ असेल, तर $PM \equal PN$.
  \item $M \equal N$ असेल, तर $MQ \equal NQ$.
  \item $M \equal N$ असेल, तर $\lambd[x][M] \equal \lambd[x][N]$.
  \item $(\lambd[x][N])Q \equal \Subst{N}{Q}{x}$.
  \end{enumerate}
\end{defn}

पहिले तीन नियम या संबंधाला तुल्यता संबंध बनवतात;
पुढील तीन त्याला रचनांशी सुसंगत बनवतात; आणि शेवटचा
नियम त्यात $\beta$-संकुचन समाविष्ट करतो.

अनौपचारिकपणे, दोन पदे $\beta$-सममूल्य असतात तेव्हाच
एका पदाचे दुसऱ्यात शून्य किंवा अधिक पायऱ्यांनी
परिवर्तन करता येते, आणि प्रत्येक पायरी $\beta$-संकुचनाची
किंवा $\beta$-संकुचनाच्या ``उलट'' प्रकारची असते.
उलट $\beta$-संकुचन
असे परिभाषित करतात की $M$ चे उलट $\beta$-संकुचन
$N$ मध्ये होते तेव्हाच $N$ चे $\beta$-संकुचन $M$
मध्ये होते.

वरील नियमांखेरीज आणखी नियम जोडून हा संबंध विस्तारू;
विस्तारित सममूल्यता संबंध $\equal[X]$ असा दर्शवू,
जेथे $X$ हा जोडलेला नियम आहे.

\end{document}
